\section{Which analytic hypothesis is actually used?}
\label{sec:complexity-grh}

In this section \(A_{\rm Dir}\) denotes precisely the supplied theorem
\texttt{LFunction\_eq\_zero\_re} in \texttt{ALPOGE\_Main.lean}, whose type is
\texttt{DirichletLZerosOnLine} in \texttt{ALPOGE\_Statements.lean}.
It quantifies over Dirichlet characters of every positive modulus and zeros in
\(0<\Re(s)<1\).  The stronger nonvanishing statement
\texttt{LFunction\_ne\_zero\_of\_half\_lt\_re} has type
\texttt{DirichletLNonvanishingRightOfHalf}.  Neither type quantifies over number
fields, their Dedekind zeta functions, or Hecke characters of such fields.
The following consequences use the supplied theorems as cited premises; they
do not assert that the missing imported Lean modules were rebuilt in this run.
The cited premise is Alp\"oge's supplied note and formal statements
\cite{CAlpoge26}.

\subsection{The part that passes to number fields}

\begin{proposition}[C-GRH-01: the cyclotomic-field transfer]
Assume \(A_{\rm Dir}\).  Let \(K\) be a subfield of a cyclotomic field
\(\mathbb Q(\zeta_N)\).  All zeros of \(\zeta_K(s)\) in
\(0<\Re(s)<1\) lie on \(\Re(s)=1/2\).
\end{proposition}
\begin{proof}
Put \(G=\operatorname{Gal}(K/\mathbb Q)\), a finite abelian group, and let
\(X_K\) be its character group, identified with the primitive Dirichlet
characters inducing the characters of this quotient of
\((\mathbb Z/N\mathbb Z)^\times\).  We verify the factorization
\[
 \zeta_K(s)=\prod_{\chi\in X_K}L(s,\chi).
\]
At a rational prime \(p\), let \(I\subseteq G\) be inertia, let
\(e=|I|\), and let the Frobenius element of \(G/I\) have order \(f\).
There are \(g=|G|/(ef)\) primes above \(p\), each of residue degree \(f\).
The Dedekind local factor is therefore \((1-p^{-fs})^{-g}\).
The Dirichlet factors that have a local factor at \(p\) are exactly those
characters trivial on \(I\).  On the cyclic subgroup generated by Frobenius,
their values run through all \(f\)-th roots of unity, each \(g\) times.
Thus their product is also \((1-p^{-fs})^{-g}\).  Equality of the absolutely
convergent Euler products proves the displayed identity for \(\Re(s)>1\),
and meromorphic continuation proves it everywhere.  In the open critical
strip each factor is holomorphic; the possible principal-character pole is
at \(s=1\), outside this strip.  A zero of their finite product must be a
zero of one factor, and \(A_{\rm Dir}\) puts that zero on the critical line.
\end{proof}

This is a genuine enlargement of the directly named family, but not an
extension to arbitrary number fields.  The standard Kronecker--Weber theorem
identifies these fields with the finite abelian extensions of \(\mathbb Q\);
the proposition was stated for cyclotomic subfields so that its proof needs
only the explicit character dictionary.  This dictionary is also explained
in Bach \cite[\S6, p.~369]{CBach90}.

\subsection{Koiran's theorem and its remaining gate}

The exact result needed is Koiran's \emph{Hilbert's Nullstellensatz is in the
polynomial hierarchy}, DIMACS long version, Theorem~2: under the paper's
analytic hypothesis, \(\mathrm{HN}\in\mathrm{AM}\).  Here HN asks whether
integer polynomial equations have a common complex zero.  The analytic
ingredient is \S5, Theorem~9 and Corollary~1, for an \emph{arbitrary}
irreducible \(f\in\mathbb Z[X]\).  If \(d=\deg f\), \(\Delta\) is its
discriminant, and \(W_f(p)\) is its number of roots modulo \(p\), Theorem~9
bounds the discrepancy
\[
 \left|\sum_{\substack{p\le x\\p\nmid\Delta}}
                 (1-W_f(p))\right|
 =O\!\left(x^{1/2}\log(|\Delta|x^d)\right).
\]
Koiran explicitly says that an effective prime ideal theorem suffices.  The
field is \(\mathbb Q(\alpha)\) for a primitive algebraic solution coordinate;
it need not be abelian.  The 8 August 1997 erratum corrects the quantitative
bound in Theorem~8, without changing Theorems~1 or~2
\cite[Theorem~2; \S5, Theorem~9, Corollary~1; erratum]{CKoiran96}.

Consequently \(A_{\rm Dir}\) discharges this input for cyclotomic subfields,
by C-GRH-01, but the supplied statement does not instantiate the theorem for
general HN.  This is a \emph{scope gap}, not a proof of the negation of
\(A_{\rm Dir}\Rightarrow\mathrm{HN}\in\mathrm{AM}\).  The exact unresolved
gate is a uniform effective prime-ideal estimate of the above strength for
the arbitrary primitive-element fields arising in the reduction.  One may
use a weaker suitable analytic statement instead of full Dedekind GRH:
Rojas \cite[Theorem~1 and Theorem~2]{CRojas03} gives the chain
\[
 \text{Dedekind GRH}\Longrightarrow\mathrm{DZH}
 \Longrightarrow\mathrm{GIPIT}\Longrightarrow\mathrm{HN}\in\mathrm{AM}.
\]
Neither DZH nor GIPIT for these arbitrary fields is supplied by a theorem
about Dirichlet characters alone.

\begin{proposition}[C-GRH-03: an HN instance outside the abelian-field restriction]
The feasible Hilbert--Nullstellensatz instance \(X^3-2=0\) has no
solution in any finite abelian extension of \(\mathbb Q\).
Consequently, feasible HN instances cannot all be handled by restricting
their solution fields to those covered by C-GRH-01.
\end{proposition}
\begin{proof}
The polynomial has a real root, so the instance is feasible.
Eisenstein's criterion at \(2\) proves that it is irreducible over
\(\mathbb Q\).  For any of its roots \(\alpha\), therefore,
\([\mathbb Q(\alpha):\mathbb Q]=3\).
If this degree-three extension were Galois, the irreducible polynomial
would split in it.  Taking the ratio of two distinct roots would put
a primitive cube root of unity \(\zeta_3\) in \(\mathbb Q(\alpha)\).
But \([\mathbb Q(\zeta_3):\mathbb Q]=2\), contradicting the tower law.
Thus no root field is Galois.  Every intermediate field of a finite
abelian Galois extension of \(\mathbb Q\) is Galois, since every
subgroup of an abelian group is normal.  No root \(\alpha\) can belong
to such an extension, proving the assertion.
\end{proof}

This example excludes the proposed restriction of solution fields;
it does not prove that Dirichlet GRH cannot have other consequences
for HN.

\subsection{Noether normalization versus identity testing}

GCT~V uses h.s.o.p. for a homogeneous system of parameters, s.s.o.p.
for a small system of parameters, and e.s.o.p. for an explicit system
of parameters; explicitness includes polynomial-time construction
\cite[\S1.2]{CGCTV}.
There are two separate inputs in GCT~V.  Its Theorem~2.10 invokes Koiran.
This is used for the conditional PH bound for constructing an h.s.o.p.
(Theorem~3.4), and for the conditional EXPH bound for constructing or
verifying an s.s.o.p. of the succinct determinant variety
(Theorem~4.10).  These inherit the number-field gate just described.
Theorem~4.1, concerning an h.s.o.p. for that succinct variety, has the further
explicit hypothesis that its defining equations can be computed in
exponential time.  Even full GRH would not erase that additional hypothesis
\cite[Theorems~2.10, 3.4, 4.1, 4.10]{CGCTV}.

The derandomization results use different hypotheses.  Theorem~4.9(a) states
that a polynomial-time construction of an e.s.o.p. for
\(\Delta[\det,m]\) is equivalent to \emph{strengthened} black-box
derandomization of symbolic determinant identity testing.  Strengthened
means a hitting set that also hits nonzero limits of the prescribed
small-circuit family.  Theorems~1.7 and 4.7 obtain a quasi-polynomial
construction from subexponential \emph{border} determinantal lower bounds
for the permanent.  Theorem~5.11 gives the corresponding statements for
general explicit varieties, using strengthened low-degree PIT or an
explicit integral multilinear family with the stated arithmetic border
lower bound.  No GRH hypothesis substitutes for these circuit or hitting-set
hypotheses.  Conversely, these equivalences and reductions are already
proved without GRH.  The remaining gate is the specified hitting set or
border lower bound, not an unproved occurrence of the word ``GRH''
\cite[Theorems~1.7, 4.7, 4.9, 5.11]{CGCTV}.

\subsection{Miller and Bach: the direct discharge}

\begin{proposition}[C-GRH-02: the witness bounds under the supplied theorem]
Assume \(A_{\rm Dir}\) and the supplied Riemann-hypothesis theorem.
The conditional hypotheses in Miller's deterministic primality test and in
Bach's rational-field bounds are discharged.  In particular an odd
composite integer \(m\) has a Miller compositeness witness at most
\(2(\ln m)^2\).  The integers below \(3(\ln m)^2\) that are coprime to
\(m\) generate \((\mathbb Z/m\mathbb Z)^\times\).
\end{proposition}
\begin{proof}
Bach's \S6 explicitly restricts its analysis to Dirichlet characters and
the Riemann zeta function.  Theorem~2, p.~372, bounds the first positive
integer outside any proper subgroup, allowing a nonunit, by
\(2(\ln m)^2\).  Theorem~3 bounds the first \emph{unit} outside any
proper subgroup by \(3(\ln m)^2\).  The consequences for witnesses and
generators are stated on pp.~372--373.  Their analytic assumptions are
exactly covered by the supplied theorems, so these cited results apply.
For the generator assertion specifically, the subgroup generated by the
specified units cannot be proper: Theorem~3 would supply a unit outside
it already among its generators, a contradiction.
\end{proof}

Miller's original polynomial-time conclusion is Theorem~2, with the
analytic hypothesis defined in his Appendix \cite{CMiller76}; Lenstra's
simplified criterion needs only a certain family of quadratic Dirichlet
characters \cite[Theorem~1]{CLenstra79}.  Ordinary deterministic
polynomial-time primality testing was already unconditional by AKS
\cite[Theorems~4.1 and 5.1]{CAKS04}.  What the supplied theorem adds here
is the validity of these specific small-witness guarantees.  Bach's
Theorem~4 concerns Hecke characters over general number fields and retains
its broader analytic gate; it must not be bundled with Theorems~2 and~3.

\section{The supplied counting theorem: quantitative consequences and limits}
\label{sec:complexity-eth}

Write \(A_{\#\mathrm{SAT}}\) for OpenAI's supplied Theorem~1.1 in
\path{LIT_OpenAI_ETH_is_False/main_full_version.tex}.  Its precise
bound is, for each fixed \(k\ge2\),
\[
 (n+m)^{O_k(1)}2^{O_k(n\log\log n/\sqrt{\log n})},
\]
with polynomial space.  For every fixed \(D\), when \(m\le Dn\), the
exponent improves to \(O_{k,D}(n/\sqrt{\log n})\).  The source is
\cite[Theorem~1.1; Appendix~D]{COAI26}.  These are bit-operation bounds,
not unit-cost bounds for symbolic arithmetic.

\subsection{ETH, SETH, counting, and the direction of reductions}

\begin{proposition}[C-SAT-01: the quantifiers that refute SETH]
\(A_{\#\mathrm{SAT}}\) refutes the usual deterministic ETH, \(\#\mathrm{ETH}\),
SETH, and their bounded-error randomized versions.  It also refutes the
usual Gap-ETH and NSETH.  A uniform bound in growing \(k\) is not needed
for this statement.
\end{proposition}
\begin{proof}
For fixed \(k\), the exponent divided by \(n\) tends to zero because
\(\log\log n/\sqrt{\log n}\to0\).  Testing whether the exact count is
positive decides satisfiability, and testing whether it is zero decides
unsatisfiability.  Taking \(k=3\) contradicts ETH and \(\#\mathrm{ETH}\).
Fix, for example, \(\varepsilon=1/2\).  For every fixed \(k\) its counting
algorithm, with its own constants and threshold, eventually runs in
\(2^{(1-\varepsilon)n}(n+m)^{O_k(1)}\) time; the finitely many smaller
lengths can be absorbed into its constant.  This is precisely the
negation of the SETH quantifier
\(\forall\varepsilon>0\ \exists k\) for which no such algorithm exists.
A deterministic algorithm is also a bounded-error randomized algorithm,
and also a nondeterministic one.  Its unsatisfiability version therefore
contradicts NSETH.  On the Gap-ETH promise, satisfiability distinguishes
the two promised cases, giving the last assertion.
\end{proof}

The conclusions of C-SAT-01 are already stated in the supplied paper,
\S1.1.1--1.1.2; the short proof above records the quantifier issue explicitly.
An implication \(\mathrm{ETH}\Rightarrow H\) remains a true conditional
theorem when ETH is false.  It no longer establishes \(H\), but it also
does not establish \(\neg H\).  The same applies to ETH-based
\(f(k)n^{o(k)}\) exclusions and to SETH-based fine-grained lower bounds.
For example, refuting SETH does not refute the Orthogonal Vectors
Conjecture merely because SETH implies that conjecture.  Nor does a
fixed-width guarantee settle the stronger hypothesis about unrestricted
CNF-SAT with width growing with \(n\).  These distinctions are already
made in \cite[\S1.1.3]{COAI26}.

\begin{proposition}[C-COUNT-01: a precise counting transfer]
Let \(V_x\) be a given Boolean verifier with \(w\) free input bits and
\(s\) gates of fan-in at most two, with all other inputs fixed by \(x\).
Under \(A_{\#\mathrm{SAT}}\), its number of accepting inputs can be computed
in
\[
 (|x|+w+s)^{O(1)}
 2^{O((w+s+1)/\sqrt{\log(w+s+2)})}
\]
bit operations and polynomial space.  The constants are absolute for
the fixed gate basis.  The construction is parsimonious.
\end{proposition}
\begin{proof}
Introduce one variable for each free input and one for each gate value.
For a two-input gate computing \(z=g(a,b)\), there are four forbidden
triples \((a,b,z)\), one for each choice of \((a,b)\).  For each forbidden
triple, insert the clause that is false precisely at that triple.
These clauses have width at most three.  A unary or constant gate has
at most four such clauses too, after substituting fixed bits.  Repeated
inputs simply produce shorter clauses or tautologies.  Pin the output
to one with a unit clause.  The formula has \(N=w+s\) variables and
at most \(4s+1\le5(N+1)\) clauses.  Once its free inputs are assigned,
the circuit's topological order determines all gate variables uniquely.
The clauses hold exactly for that assignment to gate values, and the
output condition holds exactly for accepting inputs.  This proves
parsimony.  The bounded exceptional cases \(N<20\) are enumerated.
Otherwise the fixed-density part of the supplied theorem, with \(k=3\)
and a fixed density bound, proves the stated time and space estimates.
\end{proof}

If a representation-theoretic coefficient already has a positive counting
formula with a verifier of \(B\) Boolean inputs and gates, C-COUNT-01
computes it in \(\operatorname{poly}(B)2^{O(B/\sqrt{\log B})}\).
The existence of the formula is an explicit hypothesis: this observation
does not construct a \(\#\mathrm P\) interpretation for plethysm or
Kronecker coefficients.  If \(B=O(L^a)\) in an input parameter \(L\),
the certified exponent is \(O(L^a/\sqrt{\log L})\), not automatically
\(o(L)\), and not \(O(\log L)\).  Thus polynomial reductions preserve
polynomial-time membership but need not preserve subexponential time
in the original parameter.  A target-to-SAT reduction is required to
obtain an algorithm for that target; the SAT-to-target reductions used
to prove hardness point in the other direction.

\begin{corollary}[C-COUNT-03: the parsimonious completeness form]
Counting satisfying assignments of width-at-most-three CNF formulas is
\(\#\mathrm P\)-complete under polynomial-time parsimonious reductions.
This statement does not require \(A_{\#\mathrm{SAT}}\).
\end{corollary}
\begin{proof}
Use the witness definition of \(\#\mathrm P\): a function counts the
strings \(y\in\{0,1\}^{p(|x|)}\) accepted by a fixed polynomial-time
verifier on \((x,y)\).  For the equivalent accepting-computation
definition, pad the unused nondeterministic choices by zeros and have
the verifier reject nonzero padding; this preserves the count.
Let \(T\) bound the verifier's running time and input length, and freeze a halting
configuration thereafter.  A standard one-tape simulation may enlarge
\(T\) polynomially.  Encode each tape cell by a fixed number of bits
for its symbol, head marker and finite control state.  In each time
step a cell's new data depend only on it and its two neighbors.
The finitely many local update rules have constant-size fan-in-two
Boolean circuits.  A tableau with \(T\) layers of \(O(T)\) cells
therefore gives an \(O(T^2)\)-gate circuit, constructed in polynomial
time, with \(x\) fixed and the witness bits free.  An OR tree over
the final cells tests the accepting state.  Starting configurations
are valid, so behavior of the local rules on invalid configurations
is irrelevant.  The clauses in C-COUNT-01 now give a width-three
formula with exactly one satisfying extension per accepting witness.
Conversely a satisfying assignment is itself a polynomial-length,
polynomially checkable witness.  This proves both hardness and membership.
\end{proof}

\begin{corollary}[C-SIZE-01: a concrete width and size accounting]
Under \(A_{\#\mathrm{SAT}}\), a CNF with \(n\) declared variables and
\(L\) literal occurrences can be counted in
\(\operatorname{poly}(n+L)2^{O((n+L)/\sqrt{\log(n+L+2)})}\), using
polynomial space, regardless of its width.  In particular this is
\(2^{o(n)}\operatorname{poly}(n+L)\) when
\(L=O(n(\log n)^a)\) for a fixed \(0\le a<1/2\).
\end{corollary}
\begin{proof}
Use binary OR trees for clauses, a binary AND tree for their conjunction,
and at most one negation gate per negative occurrence.  The resulting
verifier has \(n\) free inputs and \(O(L+1)\) gates; empty clauses and
empty conjunctions are handled directly.  C-COUNT-01 gives the first
bound.  Under the additional assumption,
\(\log(n+L+2)=\Theta(\log(n+2))\) and the exponent is
\(O(n/(\log n)^{1/2-a})=o(n)\), proving the second assertion.
\end{proof}

This also shows concretely why the cited theorem alone does not give the
standard uniform variable-width strong-SAT conclusion: general instances
need not satisfy this incidence bound.

\subsection{A signed, weighted permanent consequence}

The supplied Boolean contraction theorem is stronger than the decision
consequence alone.  Its exact interface is Theorem~D.1 in
\texttt{quantitative-reduction.tex}: the coefficient ring is
\((\mathbb Z/2^{P_0}\mathbb Z)[V]/(V^\ell)\), with \(P_0,\ell=O(s)\);
there are \(O(s)\) Boolean wires; maximum degree is at most
\(\lceil s^{9/10}\rceil\); and local entries have the stated uniform
\(\operatorname{poly}(s+\mathcal L)2^{O(s^{9/10})}\) time and polynomial
space descriptions.  The conclusion is
\(\operatorname{poly}(s+\mathcal L)2^{O(s/\sqrt{\log s})}\) time and
polynomial space.  We denote this cited result by \(A_{\rm contract}\).

\begin{proposition}[C-COUNT-02: signed sparse permanent]
Assume \(A_{\rm contract}\).  Let \(A\) be an \(n\times n\) integer
matrix with \(E\) nonzero entries, each of absolute value at most \(2^b\),
where \(b\ge0\).  Put \(S=\max\{2,E+n(b+1)\}\).  Its exact permanent
can be computed deterministically in
\[
 (S+\mathcal L)^{O(1)}2^{O(S/\sqrt{\log S})}
\]
time and polynomial space, where \(\mathcal L\) is its input-description
length.  In particular, bounded integer weights and \(E=O(n)\) give
\(2^{O(n/\sqrt{\log n})}\) time up to a polynomial factor.
\end{proposition}
\begin{proof}
Return zero if a row or column has no nonzero entry; handle \(n=0\)
by the empty permanent, equal to one.  Otherwise introduce a Boolean
selection variable for each supported edge of the row-column bipartite
graph.  At each row and each column impose that exactly one incident
edge is selected.  Realize such a constraint of degree \(d\) by a chain
of \(d\) ternary tensors
\[
 T_j(u,x,v)=\begin{cases}1&u+x=v\text{ as integers},\\0&\text{otherwise},\end{cases}
 \qquad u,x,v\in\{0,1\},
\]
with initial state pinned to zero and final state pinned to one.
The state after each step is the number of selections so far.  There
is a unique state assignment for a selection with exactly one edge,
and no state assignment otherwise.  On a supported edge insert a
binary diagonal tensor with values \(1\) on \((0,0)\), \(a_{ij}\)
on \((1,1)\), and zero on the two unequal pairs.
The contraction is now precisely \(\operatorname{perm}(A)\): every
perfect matching contributes the product of its weights once.

All vertices have degree at most three, and the graph has \(O(E+n)\)
edges and vertices.  Subdivide every edge by a binary equality tensor.
The old and the new vertices are its two bipartition classes; equal
bits extend uniquely across a subdivision, so the contraction is
unchanged.  The size remains \(O(E+n)\).

Let \(d_i\) be the number of supported entries in row \(i\).  The
number of perfect matchings is at most \(\prod_i d_i\le2^E\), since
\(d_i\le2^{d_i}\) and \(\sum_i d_i=E\).  Each matching weight has
absolute value at most \(2^{nb}\), whence
\(|\operatorname{perm}(A)|\le2^{E+nb}\).
Use modulus \(2^P\) with \(P=E+nb+2=O(S)\), and take \(\ell=1\).
The local tensors can be evaluated in polynomial time and space,
including their integer residues.  They satisfy every hypothesis of
\(A_{\rm contract}\) with size parameter \(S\); the bounded small
exceptions to the degree inequality can be enumerated directly.
After contraction, choose the representative in
\([-2^{P-1},2^{P-1})\).  The displayed absolute bound is at most
\(2^{P-2}\), so this representative is the exact signed permanent.
The stated costs follow from the contraction theorem.
\end{proof}

For dense bounded-entry matrices this construction has \(S=\Theta(n^2)\).
It certifies an exponent \(O(n^2/\sqrt{\log n})\), which is larger
than the \(O(n)\) exponent given by the \(2^{n-1}\)-term formula in
Landsberg \cite[\S8.2, equation~(22)]{CLandsberg13}.
It therefore supplies no improved dense-permanent bound.
The signed weighted statement is an extension of the unweighted
perfect-matching case already given in \cite[Corollary~1.3]{COAI26}.
It concerns evaluation on integer inputs, not polynomial-size
arithmetic circuits for the symbolic permanent, and not membership
or nonmembership in a determinant orbit closure.

\subsection{A bounded finite-field PIT transfer}

\begin{proposition}[C-PIT-01]
Assume \(A_{\#\mathrm{SAT}}\).  Let a division-free arithmetic circuit
over a specified prime field \(\mathbb F_p\) have \(r\) inputs,
\(s\) arithmetic gates, and output polynomial of total degree \(D<p\).
Put \(b=\lceil\log_2 p\rceil\) and
\(B=\max\{2,(s+r+1)b^2\}\).  Its polynomial identity can be tested
deterministically in \(\operatorname{poly}(B)2^{O(B/\sqrt{\log B})}\)
bit operations and polynomial space.
\end{proposition}
\begin{proof}
Use \(rb\) Boolean input bits to encode a field point uniquely, rejecting
any coordinate encoding at least \(p\).  Schoolbook addition,
multiplication and modular reduction on \(b\)-bit integers have
fan-in-two Boolean circuits of \(O(b^2)\) gates.  Replacing every
arithmetic gate by these circuits, and testing whether the output is
nonzero, gives a Boolean verifier of \(O(B)\) inputs and gates.
C-COUNT-01 counts its accepting inputs within the claimed bounds.
The zero polynomial has no accepting input.  A nonzero polynomial
of degree \(D<p\) cannot vanish on all of \(\mathbb F_p^r\).
For completeness, the latter assertion follows by induction on \(r\):
write the polynomial in its last variable; if it vanishes at every
point then, at every fixed choice of the other variables, this
univariate polynomial of degree less than \(p\) has \(p\) roots,
so all its coefficients vanish.  Apply the induction hypothesis to
each coefficient, whose individual degrees are still less than \(p\).
Thus the exact count is positive if and only if the polynomial is
nonzero, completing the test.
\end{proof}

C-PIT-01 makes the model boundary explicit.  It supplies neither a
polynomial-time hitting set, nor characteristic-zero identity testing
with arbitrary coefficient growth, nor a hitting set for the
approximative families in GCT~V.  The exponent depends on Boolean
evaluation size.  Furthermore \(2^{O(B/\sqrt{\log B})}\) is not the
\(\bigcap_{\varepsilon>0}\operatorname{NTIME}(2^{B^\varepsilon})\)
bound called NSUBEXP in arithmetic hardness-versus-randomness results;
one must not substitute one convention for the other.

\subsection{What the contraction method says about positivity}

The source uses collision packing (Lemma~D.3), integral local
coordinates (Lemma~D.4), and tapered Hermite recovery (Lemma~D.5).
At depth \(d\), a network of \(M=O(s8^{-d})\) wires uses
\(u=\Theta((\log s)/2^d)\) retained valuation levels.  Weighted modes
are discarded because their coefficients are divisible by the
working power of two, and the remaining exponent counts sum as
\(O(s/\sqrt{\log s})\).  These are divisibility and interpolation
statements; the source proves no ordered-sign or Schur-positive
interpretation for them \cite[Appendix~D.4]{COAI26}.

\begin{proposition}[C-SIGN-01: divisibility is not a positivity certificate]
The existence of an \(A_{\rm contract}\)-admissible network does not
imply nonnegativity of its integer contraction.  Nor does nonnegativity
of its canonical residue modulo \(2^P\) imply nonnegativity of the
integer it represents.
\end{proposition}
\begin{proof}
Take a network with one edge, a unary tensor constantly equal to one
at one endpoint, and the unary tensor \((-1,0)\) at the other.  Its
contraction is \(-1\); its entries, degrees, size and entry costs
satisfy the theorem for every fixed \(P\), after increasing the
size parameter if necessary.  Its canonical residue is \(2^P-1\),
a positive integer.  More structurally, a nonzero finite ring cannot
carry an additive order with \(1>0\): repeated addition would give
\(0=2^P\cdot1>0\) in \(\mathbb Z/2^P\mathbb Z\).
Thus a separate positivity interpretation is indispensable.
\end{proof}

The theorem's evaluation capacity can therefore pass to a proposed
representation-theoretic counting or tensor model once that model and
its resource bounds are proved.  It does not build that model,
establish saturation, construct a canonical basis, identify an
intersection-cohomology stalk, or compare the permanent and
determinant multiplicities.  Those are the outstanding mathematical
interfaces.

\subsection{Circuit lower bounds and already considered algebraic bridges}

The supplied paper already proves the quantitative consequence
\[
 \exists L\in\mathrm E^{\mathrm{NP}}\ \exists c>0:
 C_L(n)>c n\sqrt{\log n}\quad\text{for infinitely many }n,
\]
where gates have fan-in at most two and unrestricted fan-out
\cite[Corollary~1.4; Appendix~C.1; \texttt{quantitative-circuits.tex}]{COAI26}.
The language is in \(\mathrm E^{\mathrm{NP}}\), the bound is
superlinear rather than superpolynomial, and the model is Boolean.
This supplies none of the three missing changes needed for a direct
GCT separation: an appropriate algebraic family, a sufficiently large
arithmetic lower bound, and a border/orbit-closure certificate.

A later paper already addresses this bridge.  Chukhin--Gajulapalli--Golovnev--Mihajlin,
ECCC TR26-313, Theorem~1, proves the independent disjunction
\[
 \mathrm{NE}\not\subseteq\mathrm{SIZE}[2^n/(2n)]
 \quad\text{or}\quad
 \mathrm{VP}^{0}_{\mathbb Z}\ne\mathrm{VNP}^{0}_{\mathbb Z}.
\]
Here \(0\) means constant-free circuits of polynomial formal degree,
a different model from arbitrary complex constants.  Their Theorem~6
turns sufficiently small integer arithmetic circuits for the permanent
into perfectly complete input-oblivious Merlin--Arthur protocols for
sparse 3-UNSAT, with running time \(s(Dn)n^D\), for a fixed \(D\)
and the size bound \(s\) satisfying the theorem's uniformity conditions.
Their discussion on p.~9 explicitly observes that the supplied
\(2^{O(n/\sqrt{\log n})}\) deterministic solver is an empty-message
protocol, and compares it with the stronger
\(2^{O(n/\log n)}\) protocol gate \cite[Theorems~1 and~6; pp.~8--9]{CCGGM26}.
Thus the exact remaining contrapositive gate is a lower bound against
the specified protocols at a \emph{faster} time scale.  Refuting ETH
does not establish this gate.

The same source gives a concrete positive zero-test: Theorem~5 uses
the sum of squares of coefficients of a syntactically set-multilinear
integer circuit.  With \(d\) variable blocks of size at most \(m\),
\(s\) gates, and every gate polynomial having coefficient
\(\ell_1\)-norm at most \(2^L\), its running time is
\(m^{d-\lceil d/6\rceil}(dsL\log(2m))^{O(1)}\).
This promised positive-scalar test supplies no plethysm formula
\cite[Theorem~5, \S8.2]{CCGGM26}.

\paragraph{Status of the major separations.}
None of the deductions in this section proves \(\mathrm P=\mathrm{NP}\),
\(\mathrm P\ne\mathrm{NP}\), \(\mathrm{VP}=\mathrm{VNP}\), or
\(\mathrm{VP}\ne\mathrm{VNP}\) in the characteristic-zero model of
the question.  This sentence reports the scope of the proved deductions;
it is not a claimed logical independence theorem.  In particular,
\(\exp(O(n/\sqrt{\log n}))\) is a superpolynomial \emph{bound},
so it does not certify polynomial time, and an algorithmic upper
bound does not supply a superpolynomial lower bound for the same
problem.  The remaining gates are recorded individually in
\texttt{QUESTIONS.tsv}; none is declared refuted merely because the
supplied theorem has a different type.
